% === §5.2 Estado del arte o antecedentes relevantes ===
% Redactado por el Bibliotecario (evidencia + refs) + Investigador (hipótesis al cierre).
% Estructura conforme a guiaProyectosIA_Agente.md (subrequisitos 1–5 de §5.2).
% Hipótesis final: insertada por el Orquestador desde el Investigador.

% --- Subrequisito 1: Investigar y Agrupar el Estado del Arte Global ---

El estado del arte en la intersección entre inteligencia artificial agéntica y educación en ingeniería puede agruparse en seis filosofías de abordaje, cada una con limitaciones distintas y relaciones particulares con los subproblemas SP1, SP2 y SP3 formulados en la Sección~\ref{sec:problematica}. Una revisión exhaustiva (2023--2026), soportada sobre todo en revistas indexadas en cuartiles Q1/Q2, sustenta esta clasificación.

\textbf{Grupo~A — Tutores inteligentes clásicos basados en reglas y curaduría experta.}
La primera filosofía corresponde a los sistemas tutores inteligentes (ITS) tradicionales, fundados en bases de conocimiento codificadas a mano y árboles de decisión. La revisión sistemática pionera de Zawacki-Richter y colaboradores~\cite{zawacki2019systematic} y los análisis de Holmes y Tuomi~\cite{holmes2022state} muestran que estos sistemas, aunque pedagógicamente sólidos en dominios acotados, requieren años de ingeniería del conocimiento por disciplina y son poco adaptables a las trayectorias abiertas del estudiante de ingeniería. Trabajos posteriores confirman esa rigidez: la revisión de Ouyang y Jiao~\cite{ouyang2021its} y el estudio de Salas-Pilco y Yang~\cite{salaspilco2022ai} evidencian que la mayoría mantienen arquitecturas monolíticas centradas en ejercicios cerrados, sin un \emph{modelo del aprendiz} computacional trazable. \emph{Relación con subproblemas:} este grupo deja abiertos SP1 (diagnósticos multidimensionales imposibles con indicadores cerrados) y SP2 (la personalización no escala a la diversidad de trayectorias en ingeniería), condiciones que motivan el tránsito a filosofías basadas en modelos de lenguaje.

\textbf{Grupo~B — Asistentes conversacionales reactivos basados en LLMs (``wrappers'' de APIs).}
La irrupción de los grandes modelos de lenguaje generó un cuerpo extenso de estudios que evalúan ChatGPT y asistentes similares en aula. Chan y Hu~\cite{chan_hu2023perceptions} reportan, en $n=399$ estudiantes, percepciones positivas hacia el aprendizaje personalizado pero también preocupaciones éticas y de dependencia; Chan~\cite{chan2023policy} propone un marco ecológico de política universitaria para IA generativa. La revisión de Zirar~\cite{zirar2023chatgpt} sobre $25$ estudios concluye que el uso no mediado altera la evaluación y puede inhibir el aprendizaje profundo, en línea con la revisión sistemática de Zhai, Wibowo y Li~\cite{zhai2024overreliance} (1145 fuentes, 14 articles incluidos), que evidencia cómo la sobredependencia de sistemas dialogantes degrada el pensamiento crítico y la toma de decisiones. En el plano formativo, Escalante, Pack y Barrett~\cite{escalante2023aifeedback} (cuasiexperimental con $n=48$) no encuentran diferencias en Resultados de aprendizaje entre retroalimentación de ChatGPT (GPT-4) y tutor humano; Banihashem et~al.~\cite{banihashem2024feedback} ($n=74$) muestran que la retroalimentación de ChatGPT es más descriptiva mientras que la de pares identifica mejor los problemas, sugiriendo un rol complementario. Du y Alm~\cite{du2024eap} (Teoría de la Autodeterminación) y Chang et~al.~\cite{chang2023srlchatbot} (principios de diseño para chatbots basados en autorregulación) confirman que estos sistemas carecen deliberadamente de un modelo del aprendiz, memoria pedagógica de largo plazo y orquestación de estrategias didácticas. Evaluaciones de frontera, como el desempeño de ChatGPT en el USMLE~\cite{kung2023chatgpt}, demuestran capacidad explicativa pero sin trazabilidad diagnóstica del estudiante. \emph{Relación con subproblemas:} esta filosofía cubre parcialmente SP2 (asistencia puntual) pero falla en SP1 (no mantiene un modelo del aprendiz) y consume evidencia empírica débil para SP3 (muestras pequeñas, sin diseño causal sobre aprendizaje profundo).

\textbf{Grupo~C — Modelado del aprendiz basado en deep learning y knowledge tracing (KT).}
Una tercera filosofía aborda SP1 con modelos de aprendizaje profundo para trazar el estado de conocimiento. Trabajos recientes en revistas Q1/Q2 muestran avances sustanciales: Tong y Ren~\cite{tong2025deepkt} (Scientific Reports) proponen una arquitectura dual-stream con Transformer bidireccional y carga cognitiva para generar rutas de aprendizaje personalizadas (87.5\% de precisión); Fazil et~al.~\cite{fazil2024asist} (Journal of Learning Analytics) presentan ASIST, una red Attention-aware Stacked BiLSTM con AUC~0.86--0.90; Li et~al.~\cite{li2025option} integran option tracing en DKT sobre más de un millón de registros; Wang, Chen y Zhang~\cite{wang2023daktn} (AAAI) incorporan activación local sobre secuencias históricas; Su et~al.~\cite{su2023deepkt} (Frontiers in Psychology) modelan curvas de aprendizaje con CNN y cápsulas; Bai et~al.~\cite{bai2025deckt} (Entropy) introducen contraste dual-encoder bajo datos dispersos; Xia et~al.~\cite{xia2023ktcluster} (Applied Sciences) combinan clustering con redes neuronales para construir retratos del aprendiz interpretables. \emph{Limitación común:} estos modelos estiman \emph{correctitud binaria} de respuestas de opción múltiple dominio-específicas (K-12 matemáticas, inglés), no capturan las dimensiones de profundidad (motivación, argumentación técnica, resolución auténtica multimodal) en dominios de ingeniería mediados por IA. \emph{Relación con subproblemas:} directamente conectado a SP1, pero su brecha — ausencia de diagnósticos multimodales de aprendizaje \emph{profundo} — es justo la que este proyecto resuelve.

\textbf{Grupo~D — IA agéntica y arquitecturas multiagente basadas en LLMs.}
La cuarta filosofía — la más cercana al núcleo de la propuesta — desarrolla agentes autónomos basados en LLMs. Park et~al.~\cite{park2023generative} (Stanford) introducen agentes generativos con memoria-experiencia-reflexión-planificación; Xi et~al.~\cite{xi2023rise} sintetizan el marco cerebro-percepción-acción; Wang et~al.~\cite{wang2024survey} (Frontiers of Computer Science) revisan el ecosistema y lo proyectan hacia escenarios de ingeniería; Chang et~al.~\cite{chang2024llmeval} (ACM TIST) evalúan sistemáticamente los LLMs como base agéntica, incluyendo aplicaciones educativas; Yao et~al.~\cite{yao2023react} (ICLR) proponen ReAct (razonamiento+acción intercalados); Shinn et~al.~\cite{shinn2023reflexion} introducen Reflexion (refuerzo verbal por memoria episódica). En el extremo multiagente, Yang et~al.~\cite{yang2024lamas} formalizan los LLM-based Multi-Agent Systems (LaMAS); Händler~\cite{handler2023taxonomy} propone una taxonomía autonomía-alineación; Crawford et~al.~\cite{crawford2024bmw} presentan un framework industrial para automatización por colaboración multiagente; Ronanki~\cite{ronanki2025trust} introduce un marco RACI para colaboración humano-agente confiable en ingeniería de software; Abramson et~al.~\cite{abramson2024mia} (DeepMind) demuestran agentes multimodales interactivos por imitación. Una línea paralela identifica limitaciones críticas: Ji et~al.~\cite{ji2023hallucination} (ACM Computing Surveys) y Huang et~al.~\cite{huang2025hallucination} (ACM TOIS) documentan la alucinación en generación de LLMs; Liu et~al.~\cite{liu2023promptsurvey} (ACM Computing Surveys) y Min et~al.~\cite{min2023plmsurvey} (ACM Computing Surveys) sistematizan el prompting y la adaptación a tareas; Dwivedi et~al.~\cite{dwivedi2023xai} (ACM Computing Surveys) sientan bases de explicabilidad; Roitman~\cite{roitman2026guide} consolida recientemente el protocolo MCP, el protocolo A2A y los patrones de diseño agéntico. \emph{Limitación clave:} ningún sistema en este grupo integra una ontología de roles pedagógicamente fundamentada (tutor, evaluador, generador de problemas, simulador de sistemas eléctricos/electrónicos, asistente de programación), memoria pedagógica compartida y scaffolding adaptativo para aprendizaje profundo. \emph{Relación con subproblemas:} cubre la base tecnológica para SP2 pero deja sin resolver su núcleo pedagógico.

\textbf{Grupo~E — Desarrollo de software guiado por especificaciones y software asistido por IA (spec-driven, vibe coding).}
Una quinta filosofía, reciente y alineada con el énfasis curricular del proyecto en Ingeniería de Software, aborda el desarrollo asistido por IA. Piskala~\cite{piskala2026sdd} sintetiza el Spec-Driven Development (SDD) y su filiación con GitHub Spec Kit; Feng et~al.~\cite{feng2026ssde} (FSE Companion 2026) proponen Structured Spec-Driven Engineering para generación a nivel de repositorio con LLMs; la revisión de Ernst y Bavota~\cite{ernst_bavota_2024} mapea el impacto de los LLMs en la ingeniería de software moderna; Ebert et~al.~\cite{ebert_2021}, Ozkaya~\cite{ozkaya_2023} y Sommerville~\cite{sommerville_2015} aportan el marco de referencia curricular e industrial; White el~al.~\cite{white_2023vibe} introducen el concepto de vibe coding. Estos trabajos, enriquecidos por estudios de programación autoexplicativa~\cite{lin_huang_lu_2023} y de evaluación de competencias en ingeniería~\cite{ushar_barak_2024,rivadeneira_2025}, muestran que el asistente mejora la productividad pero tiende a sustituir el razonamiento del estudiante si no se instaura una intervención pedagógica deliberada. \emph{Relación con subproblemas:} relevante para SP2 (asistente de programación) y para SP3 (escenario de validación en asignaturas de ingeniería de software), pero ninguno integra agentes pedagógicamente fundados.

\textbf{Grupo~F — Validación empírica y evaluación de tecnologías educativas basadas en IA en aula.}
La sexta filosofía examina la evidencia causal del uso de IA en entornos reales. Hooshyar et~al.~\cite{hooshyar2024systematic} (systematic review) y Zawacki-Richter et~al.~\cite{zawacki2019systematic} coinciden en que la mayoría de las intervenciones se evalúan en laboratorio, con muestras pequeñas y ausencia de grupos de control. Estudios en Latinoamérica como los de Rivadeneira et~al.~\cite{rivadeneira_2025} y Sánchez-Mendiola y Carbajal-Degante~\cite{sanchez_mendiola_2023} describen pilotos puntuales sin trazabilidad TRL ni seguimiento longitudinal. Validaciones como las de Usher y Barak~\cite{ushar_barak_2024} se realizan en ciencias y no en ingeniería. En síntesis, los cuerpos revisados confirman la ausencia simultánea de (i) validación cuasiexperimental con grupo control en aulas reales de ingeniería y (ii) una ruta de despliegue que certifique TRL~6/7. \emph{Relación con subproblemas:} corresponde directamente a SP3 — la brecha que el proyecto cierra mediante un diseño cuasiexperimental pre--post y un despliegue híbrido transferible.

% --- Subrequisitos 2 y 3: Estado actual de la tecnología (punto de partida) y brechas ---

\textbf{Punto de partida del equipo investigador.}
El LabIA (Laboratorio de Inteligencia Artificial de la UNAL) y el GCPDS (categoría A1~Minciencias) parten de un avance previo documentado en la propuesta del Semillero de Ciencia de Datos e IA (Anexo~2 del expediente), donde se consolidaron prototipos funcionales en torno a: (i) aprovechamiento del \emph{Model Context Protocol} (MCP) para integración segura entre agentes, LLMs y herramientas externas~\cite{roitman2026guide}; (ii) flujos de \emph{Spec-Driven Development} con GitHub Spec-Kit y vibe coding~\cite{piskala2026sdd,feng2026ssde}; (iii) despliegue híbrido en estaciones Apple Silicon para inferencia \emph{on-device} y VPS para orquestación; y (iv) integración con LLMs pre-entrenados de frontera (Claude, Gemini) mediante CLI. Sobre ese TRL~4 previo, la presente convocatoria potenciará la base desarrollando el \emph{modelo del aprendiz} multidimensional (OE1/SP1), la \emph{arquitectura agéntica pedagógicamente fundada} (OE2/SP2) y la \emph{validación empírica con transferencia TRL~6/7} (OE3/SP3).

\textbf{Brechas tecnológicas identificadas frente al estado del arte global.}
Del contraste entre los seis grupos y la base del equipo emergen tres brechas convergentes: \textbf{(B1)} ausencia de un modelo computacional del aprendiz que distinga \emph{desempeño aparente} (texto/código generado por IA) de \emph{comprensión genuina}, integrando señales multimodales (argumentación textual, trazas de código, registros de simulación) — brecha de SP1; \textbf{(B2)} ausencia de un ecosistema agéntico con ontología de roles pedagógicamente fundamentada y memoria compartida que orqueste de forma coherente tutoría, retroalimentación adaptativa, generación de problemas auténticos, asistencia en programación, simulación de sistemas eléctricos/electrónicos y evaluación formativa — brecha de SP2; \textbf{(B3)} ausencia de evidencia cuasiexperimental causal y de un despliegue certificado TRL~6/7 en aulas reales de ingeniería en Colombia — brecha de SP3.

% --- Subrequisito 4: Posicionar la propuesta (la novedad) ---

\textbf{Novedad científico--tecnológica de la propuesta.}
Frente a (B1), la propuesta es pionera en operacionalizar el constructo de \emph{aprendizaje profundo} (Marton y Säljö~\cite{marton1976qualitative}; Biggs~\cite{biggs1999teaching}) en métricas computacionales trazables que combinan NLP sobre producciones argumentativas, analítica de código y análisis de simulación, superando los KT de correctitud binaria del Grupo~C~\cite{tong2025deepkt,fazil2024asist,li2025option,bai2025deckt}. Frente a (B2), la propuesta es la primera arquitectura multiagente basada en LLMs con roles pedagógicamente fundados (tutor, evaluador, generador de problemas, simulador de sistemas eléctricos/electrónicos, asistente de programación), orquestada bajo MCP~\cite{roitman2026guide}, con memoria pedagógica compartida y scaffolding adaptativo inspirado en la zona de desarrollo próximo de Vygotsky~\cite{vygotsky1978mind}, integrando patrones de agencia reflexiva (ReAct~\cite{yao2023react}, Reflexion~\cite{shinn2023reflexion}) y considerando mitigación de alucinación en dominios ingenieriles~\cite{ji2023hallucination,huang2025hallucination}. Frente a (B3), el proyecto entrega validación cuasiexperimental pre--post con grupo de control no equivalente en asignaturas reales de pregrado y posgrado de la UNAL, con tamaño del efecto reportado, y un despliegue híbrido reproducible que proyecta la adopción por el sector productivo colombiano hasta TRL~6/7, conexión con industria inexistente en los Grupos~E y~F. La combinación de las tres novedades en una sola cadena de valor (diagnóstico $\rightarrow$ orquestación agéntica $\rightarrow$ validación empírica transferible) constituye la contribución diferenciadora del proyecto.

% --- Subrequisito 5: Resumen e hipótesis (el cierre lo aporta el Investigador) ---

En resumen, las estrategias más relevantes del estado del arte pueden condensarse así: los ITS basados en reglas (Grupo~A) son pedagógicamente sólidos pero no escalan~\cite{holmes2022state,zawacki2019systematic,ouyang2021its}; los asistentes conversacionales basados en LLMs (Grupo~B) mejoran el acceso personalizado pero carecen de modelo del aprendiz y de orquestación deliberada~\cite{chan_hu2023perceptions,zhai2024overreliance,escalante2023aifeedback,banihashem2024feedback}; las técnicas de deep knowledge tracing (Grupo~C) trazan estados de conocimiento pero sobre correctitud binaria~\cite{tong2025deepkt,fazil2024asist,li2025option}; las arquitecturas multiagente basadas en LLMs (Grupo~D) aportan la orquestación autónoma pero sin ontología pedagógica~\cite{park2023generative,xi2023rise,wang2024survey,yang2024lamas}; el spec-driven/vibe coding (Grupo~E) potencia la productividad en software sin intervención pedagógica~\cite{piskala2026sdd,feng2026ssde,white_2023vibe}; y la validación empírica en aula (Grupo~F) sigue siendo escasa, descontextualizada de la ingeniería latinoamericana y sin trazabilidad TRL~\cite{hooshyar2024systematic,rivadeneira_2025}. El ecosistema de agentes autónomos que esta propuesta diseñará, desarrollará y evaluará, articula de manera novedosa los aportes de estos seis grupos para cerrar, de forma encadenada, las brechas B1, B2 y B3 — es decir, los subproblemas SP1, SP2 y SP3.

% [HIPOTESIS — Insertada por el Orquestador desde el Investigador]

\medskip

\noindent\textbf{La integración de un modelo computacional multidimensional del
aprendiz —que diagnostique trazablemente la motivación intrínseca, la argumentación
técnica y la resolución de problemas auténticos— con una arquitectura multiagente
pedagógicamente fundada, orquestada mediante LLMs y el Model Context Protocol (MCP),
y dotada de un mecanismo de \emph{scaffolding} adaptativo basado en la zona de
desarrollo próximo, en un ecosistema desplegado en entornos reales de aula
universitaria, produce una mejora estadísticamente significativa ($p < 0.05$) en el
aprendizaje profundo de estudiantes de ingeniería frente a la instrucción tradicional
con herramientas de IA convencionales, alcanzando un nivel de madurez tecnológica
TRL~6/7 transferible al sector productivo colombiano.}

\medskip

\noindent La hipótesis delimita con precisión las tres limitantes tecnológicas y
metodológicas que la propuesta busca resolver: (i) la ausencia de un modelo del
aprendiz que traduzca constructos educativos profundos —no reductibles a métricas de
clic o tiempo en plataforma— a representaciones computacionales trazables y
accionables por agentes autónomos; (ii) la fragmentación de las herramientas de IA
educativa, que operan sin un modelo compartido del estudiante, sin memoria pedagógica
y sin orquestación didáctica deliberada, impidiendo la personalización genuina de las
trayectorias de aprendizaje; y (iii) la carencia de evidencia causal —obtenida
mediante diseños cuasiexperimentales rigurosos en aulas reales— que respalde la
transferencia de los ecosistemas agénticos desde el laboratorio hacia el entorno
educativo y productivo. La verificación de la hipótesis se realizará mediante los
indicadores cuantitativos y cualitativos definidos en los objetivos específicos
OE1--OE3 (sección~\ref{sec:objetivos}), con el nivel de significación, el tamaño del
efecto y los criterios de verificación TRL allí establecidos.